\begin{textsample}{POS Dim 3 – 1960 – Score 11.00 – 1960/tv\_com\_1960\_103.txt}  \label{ex:f3_pos_030}
The commercial is mostly a product demonstration in a dim, wood-paneled home setting.
A \textbf{rectangular} stereo tape machine sits on a table or cabinet.
On top of it is a tall, clear-sided vertical holder filled with a neat stack of flat, light-colored tape \textbf{cartridges}, arranged one above another like a small column.
To the right, in the background, there is a table lamp with a dark shade, reinforcing the living-room atmosphere.
A hand enters from the left and reaches for the stack.
It lifts one flat \textbf{cartridge} from the top.
The \textbf{cartridge} is about the \textbf{size} of a small paperback or large index card \textbf{case}, with a \textbf{printed} label on its face.
The word “ STEREO ” is clearly visible in bold letters, along with a small picture or graphic beneath it.
The hand holds the \textbf{cartridge} above the stack, tilting it so the label faces outward.
Another hand, entering from the right with a shirt cuff visible, reaches in and takes hold of the \textbf{cartridge}, as if examining the item being explained.
The right hand separates out a single \textbf{cartridge} and moves it over the top of the vertical stack.
It lowers the \textbf{cartridge} onto the stack inside the clear guide, aligning it carefully with the others.
Once it is in place, the hand \textbf{withdraws}.
The same hand then moves down toward the control area on the front edge of the machine.
A finger presses a control or button on the front panel.
The machine ’s front is a mix of dark and light panels, with \textbf{rectangular} buttons and a small \textbf{central} control area ; the brand styling is clean and modern for the period.
The demonstration then focuses on the machine \textbf{operating} automatically.
The stacked \textbf{cartridges} remain upright in the clear holder on top while the front panel and loading area are shown below.
The hand is no longer touching the \textbf{cartridges}, emphasizing that the system is meant to handle the tapes itself.
A moment later, a hand reaches to the right side of the machine and appears to \textbf{operate} a small lever or side control.
This is followed by a view from above : the stack of “ STEREO ” \textbf{cartridges} sits on the top platform, and one \textbf{cartridge} lies partly extended into a lower tray or slot area, suggesting that a tape has been rejected or changed by the \textbf{mechanism}.
The \textbf{cartridge} label is visible as it rests in the receiving area.
The view broadens to show the whole system again in the living-room setting.
The Revere \textbf{unit} sits prominently on the tabletop, with the tall \textbf{cartridge} stack on top.
To the right of the machine are a dark-shaded lamp and small upright \textbf{rectangular} accessories, likely \textbf{speakers} or related stereo components.
A hand briefly reaches toward one of these small items on the right, then moves away, leaving the product displayed by itself.
For the final display, the tape system remains centered in the room : the main \textbf{unit} below, the \textbf{cartridge} stack above, and the lamp and accessories beside it.
Large white block lettering appears over the product reading “ new REVERE STEREO TAPE \textbf{CARTRIDGE} SYSTEM. ” In the closing image, the “ 3M COMPANY ” logo appears large at the bottom, with the Revere stereo tape \textbf{cartridge} system still visible behind the text.

% matched lemmas: cartridge, case, central, mechanism, operate, printed, rectangular, size, speaker, unit, withdraw
\end{textsample}
